\subsection{Data and Feature Engineering}
\label{sec:data}

\subsubsection{Forecasting data and spatial sampling}

The forecasting target is mean global horizontal irradiance (GHI), expressed
in W/m$^2$ at each temporal resolution. The data come from the GOES Aggregated
Physical Solar Model v4, distributed through the National Solar Radiation
Database (NSRDB) \cite{sengupta2018,nlr2026}. This satellite-informed physical
model provides gridded irradiance estimates. Factory coordinates select the
associated NSRDB grid cells; throughout the article, \emph{location} denotes
one of these grid points. The records are not on-site sensor measurements.

The ten locations form a common panel covering Brazil, Mexico, and the United
States, but the temporal tasks use two independently acquired NSRDB archives.
The official hourly archive contains native 30-min records that its pipeline
aggregates to one-hour means. The daily archive was requested separately at a
60-min interval and exported as arithmetic civil-day means; monthly GHI is then
computed as the arithmetic mean of those daily values in each civil month.
Both archives cover the same ten factory-associated locations and the nominal
2019--2024 study period, but the daily/monthly records are not reaggregations of
the exact HDF5 records used by the hourly experiment. Retrieval resolution and product revisions may therefore introduce
differences between the archives, further motivating error comparisons
within each task. Observations from 2019 provide the initial context needed
by tasks whose first forecast origin occurs in 2020.

Table~\ref{tab:variables} distinguishes model inputs from indexing and
contextual fields. TimesNet, DilatedRNN, LSTM, and XGBoost receive the same
historical GHI sequence and site identifier. Calendar variables, coordinates,
climate classes, and meteorological descriptors are excluded from their
inputs. The deterministic seasonal and climatological baselines additionally
use target calendar labels to retrieve historical values or fitted means,
as specified in the Supplementary Material.

\begin{table}[!t]
\centering
\caption{Fields used by the forecasting and contextual pipelines.}
\label{tab:variables}
\footnotesize
\setlength{\tabcolsep}{3pt}
\renewcommand{\arraystretch}{1.12}
\begin{tabular}{@{}
  >{\raggedright\arraybackslash}p{0.20\linewidth}
  >{\raggedright\arraybackslash}p{0.25\linewidth}
  >{\raggedright\arraybackslash}p{0.47\linewidth}@{}}
\toprule
\textbf{Role} & \textbf{Field} & \textbf{Use} \\
\midrule
Target and lagged input & GHI & Direct forecast target and causal history at
the task resolution. \\
Static model input & Site identifier & Embedding for neural models and one-hot
encoding for XGBoost. \\
Temporal index & Timestamp & Ordering, partition boundaries, forecast origins,
and target alignment; not an additional covariate. \\
Spatial index & Factory latitude/longitude & Selection of the NSRDB grid cell
and sampling of contextual maps; not a learned input. \\
Learned transform & Site-wise min--max scale & Fitted strictly before the
validation or test cutoff and reused without clipping. \\
Hourly physical rule & Solar elevation & Sets post-processed nighttime
forecasts to zero; not an input to model fitting. \\
Context only & K\"oppen--Geiger class and NASA POWER fields & Post-hoc
geographic and adverse-weather interpretation; never used for training,
selection, or forecasting. \\
\bottomrule
\end{tabular}
\end{table}

\subsubsection{Data quality and leakage prevention}

Each stored source is audited according to its schema before modeling.
Timestamp and GHI fields are parsed, the records are sorted chronologically,
and the resulting index must be unique, monotonic, regular, and common across
locations. Invalid timestamps, nonnumeric or non-finite GHI, negative GHI,
duplicates, discontinuous grids, and an inconsistent stored site identity are
rejected. Unsorted records are reordered before these checks. Product and
retrieval metadata are retained as provenance rather than supplied to the
forecast models. Monthly aggregation follows the daily checks.

For a site $\ell$, the matched input at origin $t$ is
\begin{equation}
 \mathbf{x}_{\ell,t}
 =
 \left(y_{\ell,t-c},\ldots,y_{\ell,t-1};\,\ell\right),
 \label{eq:matched_input}
\end{equation}
where $c$ is the task-specific context length ($C$ in
Eq.~\eqref{eq:forecast_map}) and $\ell$ is the fixed site identifier.
Every target vector starts at $t$ and must remain within its assigned
partition. Each input window contains only observations preceding its
forecast origin, while normalization parameters use only observations
before the corresponding fitting cutoff.
Hourly and daily labels identify the beginning of the target interval. Monthly records use month-end labels to identify the
civil month being forecast; these labels are not physical issuance times.
The monthly information set contains only the preceding completed months.

Min--max parameters are estimated separately for each site. Model selection
uses scales fitted to observations strictly before 1 January 2023. After the
epoch count has been selected, the refitted model uses new scales fitted
strictly before 1 January 2024. The same parameters transform the corresponding
inputs and targets and invert the forecasts. Values outside the training range
are not clipped, so extrapolation beyond that range remains visible.

\subsubsection{Geographic and climatic context}

Table~\ref{tab:locations} reports the canonical factory coordinates, the map
identifier used in Figure~\ref{fig:geographic_context}, and the
1980--2016 K\"oppen--Geiger class sampled at those coordinates. Classes and their
confidence values come from the aligned 0.008333-degree (nominal 1-km) products
of Beck et al. \cite{beck2018}. Sampling uses the raster cell containing each
WGS84 point, without interpolation. The 0.5-degree product is used only as the
background of the Americas map and never to assign a site's class.

\begin{table}[!t]
\centering
\caption{Factory reference coordinates, map identifiers, and present-day 1-km
K\"oppen--Geiger assignments.}
\label{tab:locations}
\footnotesize
\setlength{\tabcolsep}{2.5pt}
\renewcommand{\arraystretch}{1.10}
\begin{tabular}{@{}
  >{\centering\arraybackslash}p{0.07\linewidth}
  >{\raggedright\arraybackslash}p{0.31\linewidth}
  >{\raggedright\arraybackslash}p{0.12\linewidth}
  >{\raggedright\arraybackslash}p{0.24\linewidth}
  >{\centering\arraybackslash}p{0.13\linewidth}@{}}
\toprule
\textbf{Map ID} & \textbf{Factory location} & \textbf{Country} &
\textbf{Lat., lon. ($^\circ$)} & \textbf{Class (conf.)} \\
\midrule
1 & BYD Cama\c{c}ari & Brazil & $-12.6734$, $-38.2812$ & Af (100\%) \\
2 & Tesla Gigafactory Nevada & USA & $39.5404$, $-119.4391$ & BWk (100\%) \\
3 & Tesla Gigafactory Texas & USA & $30.2219$, $-97.6188$ & Cfa (100\%) \\
4 & Hyundai Metaplant Georgia & USA & $32.1605$, $-81.4510$ & Cfa (100\%) \\
5 & Rivian Normal & USA & $40.5093$, $-89.0545$ & Dfa (100\%) \\
6 & Tesla Fremont Factory & USA & $37.4927$, $-121.9415$ & Csb (100\%) \\
7 & Lucid AMP 1 Casa Grande & USA & $32.8569$, $-111.7784$ & BWh (100\%) \\
8 & GM Factory Zero & USA & $42.3815$, $-83.0453$ & Dfa (100\%) \\
9 & Ford Rouge EV Center & USA & $42.3056$, $-83.1676$ & Dfa (100\%) \\
10 & BMW San Luis Potos\'i & Mexico & $21.9683$, $-100.8492$ & BSh (50\%) \\
\bottomrule
\end{tabular}

\medskip
\parbox{\linewidth}{\scriptsize
Class descriptions from the distributed Beck et al. legend:
Af, tropical rainforest; BWh/BWk, hot/cold arid desert; BSh, hot arid steppe;
Cfa, temperate, no dry season, hot summer; Csb, temperate, dry summer, warm
summer; and Dfa, cold, no dry season, hot summer. Coordinates are
factory reference points; they select the associated NSRDB cell and are not
sensor coordinates.}
\end{table}


\begin{figure*}[!t]
\centering
\includegraphics[width=\textwidth]{mapa_contexto_geografico.png}
\caption{Geographic and climatic context of the ten factory-associated
locations. The Americas panel uses the present 0.5-degree K\"oppen--Geiger raster
for visual context; site assignments and confidence values are sampled from
the present nominal 1-km raster. Numbered markers match the Map ID column in
Table~\ref{tab:locations}. The United States and Detroit insets separate
spatially close labels. The map is descriptive and its climate fields are not
forecast-model inputs. Source: Beck et al. \cite{beck2018}.}
\label{fig:geographic_context}
\end{figure*}

Class confidence is the percentage of Beck et al.'s 12 ensemble maps that
assign the modal class to a cell \cite{beck2018}.
Nine sites have a class confidence of 100\%. BMW San Luis Potos\'i is the
exception: its BSh assignment has 50\% confidence and is therefore treated as
uncertain in interpretation. The Supplementary Material documents the climate codes,
confidence values, raster-cell indices, source metadata, and file hashes.

\subsubsection{Independent post-hoc meteorological context}

Daily NASA POWER fields for 2024 provide independent context for interpreting
forecast errors \cite{nasapowerdaily}. These gridded estimates describe
weather around each location. The variables are
corrected precipitation, cloud amount, all-sky surface shortwave irradiation,
and clear-sky surface shortwave irradiation in local solar time. The complete
table contains 366 dates for each of the ten locations.

The study defines three descriptive tags independently of forecast error:
precipitation-indicated when corrected precipitation is at least 1~mm/day,
high-cloud when cloud amount is at least 80\%, and adverse when either
condition holds. These study-specific thresholds are applied after forecasts
and errors have been calculated. NASA POWER values are excluded from model
inputs, training, validation, refitting, and model or metric selection.
POWER dates in local solar time are matched to the NSRDB calendar with fixed
local offsets, using factory coordinates rather than exact NSRDB grid-cell
centers. The resulting match provides approximate daily context, without
confirming weather at individual forecast hours or establishing the causes
of forecast errors.
